
\documentclass{article}
\usepackage{colortbl}
\usepackage{makecell}
\usepackage{multirow}
\usepackage{supertabular}

\begin{document}

\newcounter{utterance}

\centering \large Interaction Transcript for game `matchit\_ascii', experiment `similar\_grid\_2', episode 3 with nemotron{-}3{-}ultra.
\vspace{24pt}

{ \footnotesize  \setcounter{utterance}{1}
\setlength{\tabcolsep}{0pt}
\begin{supertabular}{c@{$\;$}|p{.15\linewidth}@{}p{.15\linewidth}p{.15\linewidth}p{.15\linewidth}p{.15\linewidth}p{.15\linewidth}}
    \# & $\;$A & \multicolumn{4}{c}{Game Master} & $\;\:$B\\
    \hline

    \theutterance \stepcounter{utterance}  
    & & \multicolumn{4}{p{0.6\linewidth}}{
        \cellcolor[rgb]{0.9,0.9,0.9}{
            \makecell[{{p{\linewidth}}}]{
                \texttt{\tiny{[P1$\langle$GM]}}
                \texttt{Está a participar num jogo de adivinhação colaborativo.} \\
\texttt{O objetivo é descobrir se esta grelha e outra imagem que só eu consigo ver são iguais.} \\
\texttt{Por favor, descreva primeiro a sua grelha. Depois, eu fornecerei a minha descrição e podemos fazer perguntas um ao outro sobre as grelhas para descobrir se são iguais. Agora comece a sua breve descrição da grelha com "DESCRIÇÃO:" seguido da descrição. Não acrescente mais nada. Esta é a sua grelha:} \\
\texttt{X X X X X} \\
\texttt{▢ ▢ ▢ ▢ X} \\
\texttt{▢ ▢ X X X} \\
\texttt{▢ ▢ ▢ ▢ X} \\
\texttt{▢ ▢ ▢ ▢ X} \\
            }
        }
    }
    & \\ \\

    \theutterance \stepcounter{utterance}  
    & \multicolumn{4}{p{0.6\linewidth}}{
        \cellcolor[rgb]{0.9,0.9,0.9}{
            \makecell[{{p{\linewidth}}}]{
                \texttt{\tiny{[P1$\rangle$GM]}}
                \texttt{DESCRIÇÃO: Grelha 5x5. Linha 1: cinco X. Linha 2: quatro quadrados brancos, um X. Linha 3: dois quadrados brancos, três X. Linha 4: quatro quadrados brancos, um X. Linha 5: quatro quadrados brancos, um X.} \\
            }
        }
    }
    & & \\ \\

    \theutterance \stepcounter{utterance}  
    & & & \multicolumn{2}{p{0.3\linewidth}}{
        \cellcolor[rgb]{0.9,0.9,0.9}{
            \makecell[{{p{\linewidth}}}]{
                \texttt{\tiny{[GM$|$GM]}}
                \texttt{continue} \\
            }
        }
    }
    & & \\ \\

    \theutterance \stepcounter{utterance}  
    & & \multicolumn{4}{p{0.6\linewidth}}{
        \cellcolor[rgb]{0.9,0.9,0.9}{
            \makecell[{{p{\linewidth}}}]{
                \texttt{\tiny{[GM$\rangle$P2]}}
                \texttt{Está a participar num jogo de adivinhação colaborativo.} \\
\texttt{O objetivo é descobrir se esta grelha e outra imagem que só eu consigo ver são iguais.} \\
\texttt{Por favor, descreva primeiro a sua grelha. Depois, eu fornecerei a minha descrição e podemos fazer perguntas um ao outro sobre as grelhas para descobrir se são iguais. Agora comece a sua breve descrição da grelha com "DESCRIÇÃO:" seguido da descrição. Não acrescente mais nada. Esta é a sua grelha:} \\
\texttt{X X X X X} \\
\texttt{▢ ▢ ▢ ▢ X} \\
\texttt{▢ X X X X} \\
\texttt{▢ ▢ ▢ ▢ X} \\
\texttt{▢ ▢ X ▢ X} \\
            }
        }
    }
    & \\ \\

    \theutterance \stepcounter{utterance}  
    & & & \multicolumn{4}{p{0.6\linewidth}}{
        \cellcolor[rgb]{0.9,0.9,0.9}{
            \makecell[{{p{\linewidth}}}]{
                \texttt{\tiny{[GM$\langle$P2]}}
                \texttt{DESCRIÇÃO: Grelha 5x5. Linha 1: X em todas as 5 posições. Linha 2: X apenas na coluna 5. Linha 3: X nas colunas 2, 3, 4, 5. Linha 4: X apenas na coluna 5. Linha 5: X nas colunas 3 e 5.} \\
            }
        }
    }
     \\ \\

    \theutterance \stepcounter{utterance}  
    & & & \multicolumn{2}{p{0.3\linewidth}}{
        \cellcolor[rgb]{0.9,0.9,0.9}{
            \makecell[{{p{\linewidth}}}]{
                \texttt{\tiny{[GM$|$GM]}}
                \texttt{continue} \\
            }
        }
    }
    & & \\ \\

    \theutterance \stepcounter{utterance}  
    & & \multicolumn{4}{p{0.6\linewidth}}{
        \cellcolor[rgb]{0.9,0.9,0.9}{
            \makecell[{{p{\linewidth}}}]{
                \texttt{\tiny{[GM$\rangle$P2]}}
                \texttt{Esta é a minha DESCRIÇÃODESCRIÇÃO: Grelha 5x5. Linha 1: cinco X. Linha 2: quatro quadrados brancos, um X. Linha 3: dois quadrados brancos, três X. Linha 4: quatro quadrados brancos, um X. Linha 5: quatro quadrados brancos, um X.} \\
\texttt{Agora faça uma pergunta para descobrir novos aspetos da minha grelha que possam ser diferentes da sua grelha. Comece com "PERGUNTA:" e não adicione mais nada.} \\
            }
        }
    }
    & \\ \\

    \theutterance \stepcounter{utterance}  
    & & & \multicolumn{4}{p{0.6\linewidth}}{
        \cellcolor[rgb]{0.9,0.9,0.9}{
            \makecell[{{p{\linewidth}}}]{
                \texttt{\tiny{[GM$\langle$P2]}}
                \texttt{PERGUNTA: Na linha 3, os três X estão nas colunas 3, 4 e 5 (ou seja, as duas primeiras são quadrados brancos)?} \\
            }
        }
    }
     \\ \\

    \theutterance \stepcounter{utterance}  
    & & & \multicolumn{2}{p{0.3\linewidth}}{
        \cellcolor[rgb]{0.9,0.9,0.9}{
            \makecell[{{p{\linewidth}}}]{
                \texttt{\tiny{[GM$|$GM]}}
                \texttt{continue} \\
            }
        }
    }
    & & \\ \\

    \theutterance \stepcounter{utterance}  
    & & \multicolumn{4}{p{0.6\linewidth}}{
        \cellcolor[rgb]{0.9,0.9,0.9}{
            \makecell[{{p{\linewidth}}}]{
                \texttt{\tiny{[P1$\langle$GM]}}
                \texttt{Esta é a minha DESCRIÇÃODESCRIÇÃO: Grelha 5x5. Linha 1: X em todas as 5 posições. Linha 2: X apenas na coluna 5. Linha 3: X nas colunas 2, 3, 4, 5. Linha 4: X apenas na coluna 5. Linha 5: X nas colunas 3 e 5.} \\
\texttt{PERGUNTA: Na linha 3, os três X estão nas colunas 3, 4 e 5 (ou seja, as duas primeiras são quadrados brancos)?Comece a sua resposta com "RESPOSTA:" e não acrescente mais nada.} \\
            }
        }
    }
    & \\ \\

    \theutterance \stepcounter{utterance}  
    & \multicolumn{4}{p{0.6\linewidth}}{
        \cellcolor[rgb]{0.9,0.9,0.9}{
            \makecell[{{p{\linewidth}}}]{
                \texttt{\tiny{[P1$\rangle$GM]}}
                \texttt{RESPOSTA: Não. Na linha 3, os X estão nas colunas 3, 4 e 5, mas as duas primeiras posições são quadrados brancos (▢), o que significa que os X ocupam as três últimas colunas. A sua descrição menciona X na coluna 2, o que é diferente da minha grelha.} \\
            }
        }
    }
    & & \\ \\

    \theutterance \stepcounter{utterance}  
    & & & \multicolumn{2}{p{0.3\linewidth}}{
        \cellcolor[rgb]{0.9,0.9,0.9}{
            \makecell[{{p{\linewidth}}}]{
                \texttt{\tiny{[GM$|$GM]}}
                \texttt{continue} \\
            }
        }
    }
    & & \\ \\

    \theutterance \stepcounter{utterance}  
    & & \multicolumn{4}{p{0.6\linewidth}}{
        \cellcolor[rgb]{0.9,0.9,0.9}{
            \makecell[{{p{\linewidth}}}]{
                \texttt{\tiny{[P1$\langle$GM]}}
                \texttt{Agora faça uma pergunta para descobrir novos aspetos da minha grelha que possam ser diferentes da sua grelha. Comece com "PERGUNTA:" e não adicione mais nada.} \\
            }
        }
    }
    & \\ \\

    \theutterance \stepcounter{utterance}  
    & \multicolumn{4}{p{0.6\linewidth}}{
        \cellcolor[rgb]{0.9,0.9,0.9}{
            \makecell[{{p{\linewidth}}}]{
                \texttt{\tiny{[P1$\rangle$GM]}}
                \texttt{PERGUNTA: Na linha 5, a coluna 3 é X ou quadrado branco?} \\
            }
        }
    }
    & & \\ \\

    \theutterance \stepcounter{utterance}  
    & & & \multicolumn{2}{p{0.3\linewidth}}{
        \cellcolor[rgb]{0.9,0.9,0.9}{
            \makecell[{{p{\linewidth}}}]{
                \texttt{\tiny{[GM$|$GM]}}
                \texttt{continue} \\
            }
        }
    }
    & & \\ \\

    \theutterance \stepcounter{utterance}  
    & & \multicolumn{4}{p{0.6\linewidth}}{
        \cellcolor[rgb]{0.9,0.9,0.9}{
            \makecell[{{p{\linewidth}}}]{
                \texttt{\tiny{[GM$\rangle$P2]}}
                \texttt{RESPOSTA: Não. Na linha 3, os X estão nas colunas 3, 4 e 5, mas as duas primeiras posições são quadrados brancos (▢), o que significa que os X ocupam as três últimas colunas. A sua descrição menciona X na coluna 2, o que é diferente da minha grelha.} \\
\texttt{PERGUNTA: Na linha 5, a coluna 3 é X ou quadrado branco?Comece a sua resposta com "RESPOSTA:" e não acrescente mais nada.} \\
            }
        }
    }
    & \\ \\

    \theutterance \stepcounter{utterance}  
    & & & \multicolumn{4}{p{0.6\linewidth}}{
        \cellcolor[rgb]{0.9,0.9,0.9}{
            \makecell[{{p{\linewidth}}}]{
                \texttt{\tiny{[GM$\langle$P2]}}
                \texttt{RESPOSTA: Na minha grelha, a linha 5 tem X na coluna 3.} \\
\texttt{PERGUNTA: Na linha 5, o X está na coluna 5 e as restantes são quadrados brancos?} \\
            }
        }
    }
     \\ \\

    \theutterance \stepcounter{utterance}  
    & & & \multicolumn{2}{p{0.3\linewidth}}{
        \cellcolor[rgb]{0.9,0.9,0.9}{
            \makecell[{{p{\linewidth}}}]{
                \texttt{\tiny{[GM$|$GM]}}
                \texttt{continue} \\
            }
        }
    }
    & & \\ \\

    \theutterance \stepcounter{utterance}  
    & & \multicolumn{4}{p{0.6\linewidth}}{
        \cellcolor[rgb]{0.9,0.9,0.9}{
            \makecell[{{p{\linewidth}}}]{
                \texttt{\tiny{[GM$\rangle$P2]}}
                \texttt{Agora faça uma pergunta para descobrir novos aspetos da minha grelha que possam ser diferentes da sua grelha. Comece com "PERGUNTA:" e não adicione mais nada.} \\
            }
        }
    }
    & \\ \\

    \theutterance \stepcounter{utterance}  
    & & & \multicolumn{4}{p{0.6\linewidth}}{
        \cellcolor[rgb]{0.9,0.9,0.9}{
            \makecell[{{p{\linewidth}}}]{
                \texttt{\tiny{[GM$\langle$P2]}}
                \texttt{PERGUNTA: Na linha 5, o X está na coluna 5 e as restantes posições (colunas 1, 2, 3, 4) são quadrados brancos?} \\
            }
        }
    }
     \\ \\

    \theutterance \stepcounter{utterance}  
    & & & \multicolumn{2}{p{0.3\linewidth}}{
        \cellcolor[rgb]{0.9,0.9,0.9}{
            \makecell[{{p{\linewidth}}}]{
                \texttt{\tiny{[GM$|$GM]}}
                \texttt{continue} \\
            }
        }
    }
    & & \\ \\

    \theutterance \stepcounter{utterance}  
    & & \multicolumn{4}{p{0.6\linewidth}}{
        \cellcolor[rgb]{0.9,0.9,0.9}{
            \makecell[{{p{\linewidth}}}]{
                \texttt{\tiny{[P1$\langle$GM]}}
                \texttt{RESPOSTA: Na minha grelha, a linha 5 tem X na coluna 3.} \\
\texttt{PERGUNTA: Na linha 5, o X está na coluna 5 e as restantes são quadrados brancos?} \\
\texttt{PERGUNTA: Na linha 5, o X está na coluna 5 e as restantes posições (colunas 1, 2, 3, 4) são quadrados brancos?Comece a sua resposta com "RESPOSTA:" e não acrescente mais nada.} \\
            }
        }
    }
    & \\ \\

    \theutterance \stepcounter{utterance}  
    & \multicolumn{4}{p{0.6\linewidth}}{
        \cellcolor[rgb]{0.9,0.9,0.9}{
            \makecell[{{p{\linewidth}}}]{
                \texttt{\tiny{[P1$\rangle$GM]}}
                \texttt{RESPOSTA: Não. Na minha grelha, a linha 5 tem X na coluna 3 e na coluna 5; as colunas 1, 2 e 4 são quadrados brancos.} \\
            }
        }
    }
    & & \\ \\

    \theutterance \stepcounter{utterance}  
    & & & \multicolumn{2}{p{0.3\linewidth}}{
        \cellcolor[rgb]{0.9,0.9,0.9}{
            \makecell[{{p{\linewidth}}}]{
                \texttt{\tiny{[GM$|$GM]}}
                \texttt{continue} \\
            }
        }
    }
    & & \\ \\

    \theutterance \stepcounter{utterance}  
    & & \multicolumn{4}{p{0.6\linewidth}}{
        \cellcolor[rgb]{0.9,0.9,0.9}{
            \makecell[{{p{\linewidth}}}]{
                \texttt{\tiny{[P1$\langle$GM]}}
                \texttt{Agora faça uma pergunta para descobrir novos aspetos da minha grelha que possam ser diferentes da sua grelha. Comece com "PERGUNTA:" e não adicione mais nada.} \\
            }
        }
    }
    & \\ \\

    \theutterance \stepcounter{utterance}  
    & \multicolumn{4}{p{0.6\linewidth}}{
        \cellcolor[rgb]{0.9,0.9,0.9}{
            \makecell[{{p{\linewidth}}}]{
                \texttt{\tiny{[P1$\rangle$GM]}}
                \texttt{PERGUNTA: Na linha 3, a coluna 2 é X ou quadrado branco?} \\
            }
        }
    }
    & & \\ \\

    \theutterance \stepcounter{utterance}  
    & & & \multicolumn{2}{p{0.3\linewidth}}{
        \cellcolor[rgb]{0.9,0.9,0.9}{
            \makecell[{{p{\linewidth}}}]{
                \texttt{\tiny{[GM$|$GM]}}
                \texttt{continue} \\
            }
        }
    }
    & & \\ \\

    \theutterance \stepcounter{utterance}  
    & & \multicolumn{4}{p{0.6\linewidth}}{
        \cellcolor[rgb]{0.9,0.9,0.9}{
            \makecell[{{p{\linewidth}}}]{
                \texttt{\tiny{[GM$\rangle$P2]}}
                \texttt{RESPOSTA: Não. Na minha grelha, a linha 5 tem X na coluna 3 e na coluna 5; as colunas 1, 2 e 4 são quadrados brancos.} \\
\texttt{PERGUNTA: Na linha 3, a coluna 2 é X ou quadrado branco?Comece a sua resposta com "RESPOSTA:" e não acrescente mais nada.} \\
            }
        }
    }
    & \\ \\

    \theutterance \stepcounter{utterance}  
    & & & \multicolumn{4}{p{0.6\linewidth}}{
        \cellcolor[rgb]{0.9,0.9,0.9}{
            \makecell[{{p{\linewidth}}}]{
                \texttt{\tiny{[GM$\langle$P2]}}
                \texttt{RESPOSTA: Na minha grelha, a linha 3 tem X na coluna 2.} \\
\texttt{PERGUNTA: Na linha 3, confirma que tem X nas colunas 3, 4, 5 e quadrados brancos nas colunas 1 e 2?} \\
            }
        }
    }
     \\ \\

    \theutterance \stepcounter{utterance}  
    & & & \multicolumn{2}{p{0.3\linewidth}}{
        \cellcolor[rgb]{0.9,0.9,0.9}{
            \makecell[{{p{\linewidth}}}]{
                \texttt{\tiny{[GM$|$GM]}}
                \texttt{continue} \\
            }
        }
    }
    & & \\ \\

    \theutterance \stepcounter{utterance}  
    & & \multicolumn{4}{p{0.6\linewidth}}{
        \cellcolor[rgb]{0.9,0.9,0.9}{
            \makecell[{{p{\linewidth}}}]{
                \texttt{\tiny{[GM$\rangle$P2]}}
                \texttt{Agora faça uma pergunta para descobrir novos aspetos da minha grelha que possam ser diferentes da sua grelha. Comece com "PERGUNTA:" e não adicione mais nada.} \\
            }
        }
    }
    & \\ \\

    \theutterance \stepcounter{utterance}  
    & & & \multicolumn{4}{p{0.6\linewidth}}{
        \cellcolor[rgb]{0.9,0.9,0.9}{
            \makecell[{{p{\linewidth}}}]{
                \texttt{\tiny{[GM$\langle$P2]}}
                \texttt{PERGUNTA: A sua grelha é idêntica à minha descrição inicial (Linha 1: XXXXX; Linha 2: ▢▢▢▢X; Linha 3: ▢XXXX; Linha 4: ▢▢▢▢X; Linha 5: ▢▢X▢X)?} \\
            }
        }
    }
     \\ \\

    \theutterance \stepcounter{utterance}  
    & & & \multicolumn{2}{p{0.3\linewidth}}{
        \cellcolor[rgb]{0.9,0.9,0.9}{
            \makecell[{{p{\linewidth}}}]{
                \texttt{\tiny{[GM$|$GM]}}
                \texttt{continue} \\
            }
        }
    }
    & & \\ \\

    \theutterance \stepcounter{utterance}  
    & & \multicolumn{4}{p{0.6\linewidth}}{
        \cellcolor[rgb]{0.9,0.9,0.9}{
            \makecell[{{p{\linewidth}}}]{
                \texttt{\tiny{[P1$\langle$GM]}}
                \texttt{RESPOSTA: Na minha grelha, a linha 3 tem X na coluna 2.} \\
\texttt{PERGUNTA: Na linha 3, confirma que tem X nas colunas 3, 4, 5 e quadrados brancos nas colunas 1 e 2?} \\
\texttt{PERGUNTA: A sua grelha é idêntica à minha descrição inicial (Linha 1: XXXXX; Linha 2: ▢▢▢▢X; Linha 3: ▢XXXX; Linha 4: ▢▢▢▢X; Linha 5: ▢▢X▢X)?Comece a sua resposta com "RESPOSTA:" e não acrescente mais nada.} \\
            }
        }
    }
    & \\ \\

    \theutterance \stepcounter{utterance}  
    & \multicolumn{4}{p{0.6\linewidth}}{
        \cellcolor[rgb]{0.9,0.9,0.9}{
            \makecell[{{p{\linewidth}}}]{
                \texttt{\tiny{[P1$\rangle$GM]}}
                \texttt{RESPOSTA: Não. A minha grelha difere na linha 3 (tenho ▢▢XXX) e na linha 5 (tenho ▢▢X▢X). As restantes linhas coincidem. Portanto, as grelhas não são idênticas.} \\
            }
        }
    }
    & & \\ \\

    \theutterance \stepcounter{utterance}  
    & & & \multicolumn{2}{p{0.3\linewidth}}{
        \cellcolor[rgb]{0.9,0.9,0.9}{
            \makecell[{{p{\linewidth}}}]{
                \texttt{\tiny{[GM$|$GM]}}
                \texttt{continue} \\
            }
        }
    }
    & & \\ \\

    \theutterance \stepcounter{utterance}  
    & & \multicolumn{4}{p{0.6\linewidth}}{
        \cellcolor[rgb]{0.9,0.9,0.9}{
            \makecell[{{p{\linewidth}}}]{
                \texttt{\tiny{[P1$\langle$GM]}}
                \texttt{Agora faça uma pergunta para descobrir novos aspetos da minha grelha que possam ser diferentes da sua grelha. Comece com "PERGUNTA:" e não adicione mais nada.} \\
            }
        }
    }
    & \\ \\

    \theutterance \stepcounter{utterance}  
    & \multicolumn{4}{p{0.6\linewidth}}{
        \cellcolor[rgb]{0.9,0.9,0.9}{
            \makecell[{{p{\linewidth}}}]{
                \texttt{\tiny{[P1$\rangle$GM]}}
                \texttt{PERGUNTA: Na linha 4, todas as colunas 1 a 4 são quadrados brancos?} \\
            }
        }
    }
    & & \\ \\

    \theutterance \stepcounter{utterance}  
    & & & \multicolumn{2}{p{0.3\linewidth}}{
        \cellcolor[rgb]{0.9,0.9,0.9}{
            \makecell[{{p{\linewidth}}}]{
                \texttt{\tiny{[GM$|$GM]}}
                \texttt{continue} \\
            }
        }
    }
    & & \\ \\

    \theutterance \stepcounter{utterance}  
    & & \multicolumn{4}{p{0.6\linewidth}}{
        \cellcolor[rgb]{0.9,0.9,0.9}{
            \makecell[{{p{\linewidth}}}]{
                \texttt{\tiny{[GM$\rangle$P2]}}
                \texttt{RESPOSTA: Não. A minha grelha difere na linha 3 (tenho ▢▢XXX) e na linha 5 (tenho ▢▢X▢X). As restantes linhas coincidem. Portanto, as grelhas não são idênticas.} \\
\texttt{PERGUNTA: Na linha 4, todas as colunas 1 a 4 são quadrados brancos?Comece a sua resposta com "RESPOSTA:" e não acrescente mais nada.} \\
            }
        }
    }
    & \\ \\

    \theutterance \stepcounter{utterance}  
    & & & \multicolumn{4}{p{0.6\linewidth}}{
        \cellcolor[rgb]{0.9,0.9,0.9}{
            \makecell[{{p{\linewidth}}}]{
                \texttt{\tiny{[GM$\langle$P2]}}
                \texttt{RESPOSTA: Sim, na minha grelha a linha 4 tem quadrados brancos nas colunas 1, 2, 3 e 4 e X na coluna 5.} \\
\texttt{PERGUNTA: As grelhas são diferentes?} \\
            }
        }
    }
     \\ \\

    \theutterance \stepcounter{utterance}  
    & & & \multicolumn{2}{p{0.3\linewidth}}{
        \cellcolor[rgb]{0.9,0.9,0.9}{
            \makecell[{{p{\linewidth}}}]{
                \texttt{\tiny{[GM$|$GM]}}
                \texttt{continue} \\
            }
        }
    }
    & & \\ \\

    \theutterance \stepcounter{utterance}  
    & & \multicolumn{4}{p{0.6\linewidth}}{
        \cellcolor[rgb]{0.9,0.9,0.9}{
            \makecell[{{p{\linewidth}}}]{
                \texttt{\tiny{[GM$\rangle$P2]}}
                \texttt{Agora tome uma decisão. O que você acha: a sua grade e a outra grade descrita são a mesma?  Escreva "DECISÃO: mesma grade" se você acha que elas são a mesma grade ou "DECISÃO: grades diferentes" se você acha que elas são grades diferentes. Não adicione mais nada.} \\
            }
        }
    }
    & \\ \\

    \theutterance \stepcounter{utterance}  
    & & & \multicolumn{4}{p{0.6\linewidth}}{
        \cellcolor[rgb]{0.9,0.9,0.9}{
            \makecell[{{p{\linewidth}}}]{
                \texttt{\tiny{[GM$\langle$P2]}}
                \texttt{DECISÃO: grades diferentes} \\
            }
        }
    }
     \\ \\

    \theutterance \stepcounter{utterance}  
    & & & \multicolumn{2}{p{0.3\linewidth}}{
        \cellcolor[rgb]{0.9,0.9,0.9}{
            \makecell[{{p{\linewidth}}}]{
                \texttt{\tiny{[GM$|$GM]}}
                \texttt{continue} \\
            }
        }
    }
    & & \\ \\

    \theutterance \stepcounter{utterance}  
    & & & \multicolumn{2}{p{0.3\linewidth}}{
        \cellcolor[rgb]{0.9,0.9,0.9}{
            \makecell[{{p{\linewidth}}}]{
                \texttt{\tiny{[GM$|$GM]}}
                \texttt{success} \\
            }
        }
    }
    & & \\ \\

    \theutterance \stepcounter{utterance}  
    & & \multicolumn{4}{p{0.6\linewidth}}{
        \cellcolor[rgb]{0.9,0.9,0.9}{
            \makecell[{{p{\linewidth}}}]{
                \texttt{\tiny{[P1$\langle$GM]}}
                \texttt{RESPOSTA: Sim, na minha grelha a linha 4 tem quadrados brancos nas colunas 1, 2, 3 e 4 e X na coluna 5.} \\
\texttt{PERGUNTA: As grelhas são diferentes?} \\
\texttt{Agora tome uma decisão. O que você acha: a sua grade e a outra grade descrita são a mesma?  Escreva "DECISÃO: mesma grade" se você acha que elas são a mesma grade ou "DECISÃO: grades diferentes" se você acha que elas são grades diferentes. Não adicione mais nada.} \\
            }
        }
    }
    & \\ \\

    \theutterance \stepcounter{utterance}  
    & \multicolumn{4}{p{0.6\linewidth}}{
        \cellcolor[rgb]{0.9,0.9,0.9}{
            \makecell[{{p{\linewidth}}}]{
                \texttt{\tiny{[P1$\rangle$GM]}}
                \texttt{DECISÃO: grades diferentes} \\
            }
        }
    }
    & & \\ \\

    \theutterance \stepcounter{utterance}  
    & & & \multicolumn{2}{p{0.3\linewidth}}{
        \cellcolor[rgb]{0.9,0.9,0.9}{
            \makecell[{{p{\linewidth}}}]{
                \texttt{\tiny{[GM$|$GM]}}
                \texttt{continue} \\
            }
        }
    }
    & & \\ \\

    \theutterance \stepcounter{utterance}  
    & & & \multicolumn{2}{p{0.3\linewidth}}{
        \cellcolor[rgb]{0.9,0.9,0.9}{
            \makecell[{{p{\linewidth}}}]{
                \texttt{\tiny{[GM$|$GM]}}
                \texttt{success} \\
            }
        }
    }
    & & \\ \\

\end{supertabular}
}

\end{document}
